\documentclass[11pt]{article}
% common.tex — minimal noise setup
% ----------------------------

% Silence package messages before anything else
\RequirePackage{silence}

% Completely hide "info" messages from LaTeX packages
\WarningFilter*{latex}{}
\WarningFilter*{hyperref}{}
\WarningFilter*{pgf}{}
\WarningFilter*{tikz}{}

% Optionally: hide some hyperref PDF string warnings that are common but harmless
\WarningFilter{hyperref}{Token not allowed in a PDF string}

% ----------------------------
% Now load the usual packages
\usepackage[utf8]{inputenc}
\usepackage[T1]{fontenc}
\usepackage{amsmath, amssymb, amsthm, amsfonts}
\usepackage{geometry}
\usepackage{hyperref}
\usepackage{graphicx}
\usepackage{physics}
\usepackage{booktabs}
\usepackage{listings}
\usepackage{bookmark}
\usepackage{amscd}
\usepackage{tikz-cd}
\usepackage{mathtools}
\usepackage{xspace}
\usepackage[backend=biber, style=numeric]{biblatex}
\usepackage{glossaries}
\usepackage{tcolorbox}

% ----------------------------
% Theorems, macros
\newtheorem{lemma}{Lemma}
\newcommand{\IaMe}{\(IaM^e\)\xspace}
\newcommand{\pdflink}[1]{}

% ----------------------------
% Page layout
\geometry{letterpaper, margin=1in}
\setlength{\parindent}{0pt}
\setlength{\parskip}{1em}

% Hyperref settings
\hypersetup{
    colorlinks=true,
    linkcolor=black,
    citecolor=black,
    filecolor=magenta,
    urlcolor=blue
}

% TOC formatting
\usepackage{tocloft}
\setlength{\cftbeforechapskip}{1.0em}
\setlength{\cftbeforesecskip}{0pt}
\setlength{\cftbeforesubsecskip}{0pt}
\renewcommand{\cftchapleader}{\cftdotfill{\cftdotsep}}

% Glossary
\defglsentryfmt{\ifglsused{\glslabel}{\glsgenentryfmt}{\textit{\glsgenentryfmt}}}

% Author
\author{$IaM^e$}


\input{../glossary.tex}
\addbibresource{../references.bib} 

\title{String Theory as Spectral Compression}
\author{Juha Meskanen}
\date{2025}

\begin{document}

\maketitle

\begin{abstract}
The objective is to study whether adding antisymmetric structure actually reduce description length under complexity measure.
\end{abstract}


\tableofcontents
\bigskip

\section{Hypothesis}

Compression Advantage of Antisymmetry.

For ensembles that support persistent localized structures, encodings with antisymmetric constraints achieve lower total description length than purely symmetric (bosonic) encodings.


\section{Define the Competing Models}

Define two strictly comparable representations

\subsection*{Bosonic}

\begin{itemize}
\item State = multiset of modes
\item No occupancy restriction
\item Fully symmetric under exchange
\end{itemize}

Formally:

\begin{itemize}
\item coefficients \( c_k \in \mathbb{C} \)
\item no constraint on reuse of modes
\end{itemize}

\subsection{Fermionic / Antisymmetric}

\begin{itemize}
\item State = antisymmetric combination of modes
\item Enforced uniqueness / exclusion
\item Order-sensitive with sign flips
\end{itemize}

Implementation-wise, we don’t need full QFT. We can approximate via bitstrings with no duplicate occupancy, or determinant-like encoding (Slater-style).

This corresponds to the structure behind the Pauli exclusion principle, but we’re testing it purely as an information constraint.


\section{Define Description Length Operationally}

Define a computable proxy for:
\[
L(\text{state})
\]


\section{Experimental Design (Minimal Viable Spike)}

TODO

\section{Code}

TODO

\printbibliography

\end{document}
